\textbf{Example output:}
\texttt{[``Aspirin inhibits COX-1 enzymes'',
``COX-1 inhibition reduces platelet aggregation'',
``Reduced platelet aggregation prevents blood clots'']}

\section{Proof of Theorem~\ref{thm:consistency}}
\label{app:proofB}

\textbf{Step 1 (Identifiability).}
Let the population equations for crowd-sourcing reliability
be $\Theta^* = \arg\min \mathbb{E}[\ell(\Theta, \Sigma)]$.
Under distinct marginal error rates ($p_i \neq p_j$),
the system is identified by the Dawid-Skene rank
condition~\cite{zhang2014}: the matrix of pairwise
agreement probabilities has full rank, ensuring a unique
solution for $\Theta^*$.

\textbf{Step 2 (Copula MLE consistency).}
Given identifiable $\Theta^*$, the copula log-likelihood
$\ell(\Sigma) = \log \Phi_\Sigma(\mathbf{z})$ is strictly concave
in $\Sigma$ on the bounded set $\|\Sigma\|_2 \leq B$.
The MLE $\hat{\Sigma} \xrightarrow{p} \Sigma^*$ by standard
asymptotic MLE theory~\cite{nelsen2006}.

\textbf{Step 3 (Posterior consistency).}
Given $\hat{\Theta} \xrightarrow{p} \Theta^*$ and
$\hat{\Sigma} \xrightarrow{p} \Sigma^*$,
consistency of $\hat{P}(y_c = 1)$ follows from the
continuous mapping theorem applied to the copula posterior,
which is a continuous function of $(\hat{\Theta}, \hat{\Sigma})$.
\qed

\section{Proof of Theorem~\ref{thm:guarantee}}
\label{app:proofC}

Define the non-conformity score
$s(c) = 1 - \hat{P}(y_c = 1)$.
The hallucination risk is:
$R(\lambda) = \frac{1}{N}\sum_{k=1}^N \mathbf{1}[s(c_k) \leq 1-\lambda, y_k = 0]$,
which is monotonically non-increasing in $\lambda$.

By Theorem~1 of~\citet{angelopoulos2022crc}, for any
exchangeable sequence $\{(c_k, y_k)\}_{k=1}^{N+1}$,
the threshold selected by Eq.~\eqref{eq:crc} satisfies
$\mathbb{E}[R(\lambda^*)] \leq \delta$.
The finite-sample correction
$\hat{\delta}_\alpha = \delta - \sqrt{\frac{\ln(1/\alpha)}{2N}}$
converts this expectation bound to a
high-probability bound: $P(R(\lambda^*) \leq \delta) \geq 1 - \alpha$,
which is exactly Eq.~\eqref{eq:guarantee}.

Exchangeability is guaranteed by the i.i.d.\ assumption
on calibration data and random splitting. \qed

\section{Hyperparameter Settings}
\label{app:hyperparams}

\begin{table}[h]
\caption{Full hyperparameter configuration for reproducibility.}
\label{tab:hyperparams}
\centering
\begin{tabular}{ll}
\toprule
Parameter & Value \\
\midrule
$\tau_\text{NLI}$ (entailment threshold)    & 0.85 \\
$\alpha_\text{ema}$ (EMA decay)             & 0.05 \\
$\lambda_\text{LW}$ (shrinkage intensity)   & analytical (Ledoit-Wolf) \\
$\delta$ (hallucination tolerance)          & 0.05 \\
$\alpha$ (CRC confidence complement)        & 0.10 \\
$\tau_\text{rej}$ (rejection threshold)     & 0.30 \\
BP iterations (max)                         & 50 \\
BP convergence tolerance                    & $10^{-4}$ \\
SARSOP horizon                              & 10{,}000 steps \\
POMDP observation noise $\sigma$            & 0.05 \\
Calibration set size per domain             & 500 \\
Temperature scaling (GPT-4o)               & $T = 1.15$ \\
Temperature scaling (Llama-3)              & $T = 1.25$ \\
\bottomrule
\end{tabular}
\end{table}

\end{document}